
% !TEX program = xelatex
\documentclass[12pt,a4paper]{article}

% -------------------- Encoding & language --------------------
\usepackage{fontspec}
\usepackage{polyglossia}
\setmainlanguage{russian}
\setotherlanguage{english}
\defaultfontfeatures{Ligatures=TeX,Scale=MatchLowercase}
\setmainfont{Times New Roman}
\setsansfont{Arial}
\setmonofont{Courier New}

% -------------------- Math & layout --------------------
\usepackage{amsmath,amssymb,mathtools,bm}
\usepackage{physics}
\usepackage{siunitx}
\sisetup{per-mode=symbol,detect-all=true,locale=DE}
\usepackage[margin=2.1cm]{geometry}
\usepackage{hyperref}
\usepackage{cleveref}
\usepackage{enumitem}
\usepackage{xcolor}
\usepackage{booktabs}
\usepackage{microtype}
\usepackage{tikz}

\hypersetup{colorlinks=true, linkcolor=blue!60!black, urlcolor=blue!60!black, citecolor=blue!60!black}

\title{Некомпланарные минимально-импульсные переходы в линейной теории:\\
\large Полный пошаговый вывод формул из принципа максимума и линейных уравнений}
\author{Подготовлено специально для В.\,Чиняева}
\date{\today}

\newcommand{\R}{\mathbb{R}}
\newcommand{\T}{\mathrm{T}}
\newcommand{\dd}{\,\mathrm{d}}
\newcommand{\half}{\tfrac{1}{2}}
\newcommand{\vv}[1]{\bm{#1}}
\newcommand{\uv}[1]{\hat{\bm{#1}}}

\begin{document}
\maketitle
\tableofcontents

\section{Выбор переменных и масштабирование}

\paragraph{Цель.}
Перенести вывод из статьи в \textbf{пошаговую форму}: от уравнений Гаусса $\to$ линейная модель (первого порядка по $e,i$) $\to$ постановка оптимального управления $\to$ ПМП и \emph{primer}-вектор $\to$ геометрия экстремума $p(\theta)$ $\to$ двухимпульсное решение и замыкание по измеряемым приращениям $\Delta(a/a_0),\Delta\vv e,\Delta\vv i$.

\subsection{Локальные орбитальные элементы первого порядка}
Рассматривается близость к \emph{круговой} опорной орбите с большой полуосью $a_0$ и гравитационным параметром $K=\mu$.
Вводим \emph{плоские} компоненты вектора эксцентриситета и \emph{наклонный} вектор\footnote{Достаточно линейной аппроксимации: $e=\sqrt{e_x^2+e_y^2}$ и $i=\sqrt{i_x^2+i_y^2}$ малы.}:
\begin{equation}
\vv x \equiv
\begin{bmatrix}
x_1\\x_2\\x_3\\x_4\\x_5
\end{bmatrix}
=
\begin{bmatrix}
a/a_0\\ e_y\\ e_x\\ i_y\\ i_x
\end{bmatrix}.
\end{equation}
Ось $x$ направим по линии узлов опорной орбиты, $y$ --- по касательной в узле, $z$ --- по нормали к опорной плоскости.

\subsection{Управление и независимая переменная}
Пусть $\vv \alpha(t)$ --- малое (импульсное) ускорение двигателя; разложение в базисе R--T--N:
\(
\alpha_R,\alpha_T,\alpha_N
\).
Обозначим \emph{дозу тяги} $u(t)=\int_0^t \|\vv\alpha(\tau)\|\dd\tau$ и \emph{направляющие косинусы} $l_R=\alpha_R/\|\vv\alpha\|$, $l_T=\alpha_T/\|\vv\alpha\|$, $l_N=\alpha_N/\|\vv\alpha\|$.
Переходим от $t$ к $u$ (по сути, нормируем скорость изменения на модуль управления).

\section{Линеаризация уравнений движения}

\subsection{Уравнения Гаусса и их упрощение}
Классические уравнения Гаусса для $(a,e,i,\omega,\Omega)$ при воздействии $(\alpha_R,\alpha_T,\alpha_N)$ вблизи круговой орбиты $e,i\ll1$ приводятся к линейным уравнениям для $(a/a_0,e_x,e_y,i_x,i_y)$ при выборе независимой переменной $\theta$ (центрального угла).
Переход от $\theta$ к дозе $u$ даёт (все шаги развёрнуты в Доп.\,А):
\begin{align}
\frac{\dd (a/a_0)}{\dd u} &= 2\sqrt{\frac{a_0}{K}}\,l_T, \label{eq:lin_a}\\
\frac{\dd e_y}{\dd u} &= \sqrt{\frac{a_0}{K}}\,(2l_T\sin\theta - l_R\cos\theta),\label{eq:lin_ey}\\
\frac{\dd e_x}{\dd u} &= \sqrt{\frac{a_0}{K}}\,(2l_T\cos\theta + l_R\sin\theta),\label{eq:lin_ex}\\
\frac{\dd i_y}{\dd u} &= \sqrt{\frac{a_0}{K}}\,l_N\sin\theta,\label{eq:lin_iy}\\
\frac{\dd i_x}{\dd u} &= \sqrt{\frac{a_0}{K}}\,l_N\cos\theta.\label{eq:lin_ix}
\end{align}
\textbf{Комментарий по индексам.} В сканах часто путаются $x_2/x_3$ и $x_4/x_5$. Ниже строго используем
$e_y\leftrightarrow x_2$, $e_x\leftrightarrow x_3$, $i_y\leftrightarrow x_4$, $i_x\leftrightarrow x_5$.

\subsection{Матрица чувствительности импульса}
Импульс $\Delta\vv v$ в момент $\theta$ даёт изменение элементов
\begin{equation}\label{eq:Gmap}
\Delta\vv x = \sqrt{\frac{a_0}{K}}\,
\underbrace{\begin{bmatrix}
0 & 2 & 0\\
-\cos\theta & 2\sin\theta & 0\\
\ \ \sin\theta & 2\cos\theta & 0\\
0 & 0 & \sin\theta\\
0 & 0 & \cos\theta
\end{bmatrix}}_{\displaystyle \mathbf{G}(\theta)}
\Delta\vv v_{(R,T,N)}.
\end{equation}
Это следует из интегрирования \eqref{eq:lin_a}--\eqref{eq:lin_ix} на бесконечно малом интервале $u$ (см. Доп.\,А).

\section{Оптимальная задача и принцип максимума}
\subsection{Постановка}
Требуется зафиксированное $\Delta\vv x$ между начальным и конечным состояниями при минимизации суммарного импульса
\(
J=\sum_k \|\Delta\vv v_k\|
\).
При импульсной постановке минимизация $J$ эквивалентна максимизации Гамильтониана
\(
H=\vv\lambda^\T \dot{\vv x}
\)
по направлению управления, где $\vv\lambda$ --- сопряжённый вектор.

\subsection{Primer-вектор}
Подставляя \eqref{eq:lin_a}--\eqref{eq:lin_ix} в $H$ и группируя множители при $(l_R,l_T,l_N)$, получаем
\begin{align}
\lambda &\equiv \sqrt{\frac{a_0}{K}}(-\lambda_2\cos\theta + \lambda_3\sin\theta),\\
\mu     &\equiv \sqrt{\frac{a_0}{K}}(2\lambda_1 + 2\lambda_2\sin\theta + 2\lambda_3\cos\theta),\\
\nu     &\equiv \sqrt{\frac{a_0}{K}}(\lambda_4\sin\theta + \lambda_5\cos\theta).
\end{align}
Оптимальное направление импульса совпадает с направлением $\vv p=(\lambda,\mu,\nu)$, а \emph{необходимость} оптимальности указывает, что импульсы должны происходить в глобальных максимумах $p(\theta)=\|\vv p(\theta)\|$.

\subsection{Приведение к фазовому сдвигу}
Введём угол
\(
\tan\theta_0=\lambda_2/\lambda_3
\).
Тогда
\begin{align}
\lambda &= R\,\sin(\theta-\theta_0),\qquad R=\sqrt{\frac{a_0}{K}}\sqrt{\lambda_2^2+\lambda_3^2},\\
\mu &= 2\sqrt{\frac{a_0}{K}}\left(\lambda_1 + R\cos(\theta-\theta_0)\right),\\
\nу &= \sqrt{\frac{a_0}{K}}\frac{(\lambda_3\lambda_4-\lambda_2\lambda_5)\sin(\theta-\theta_0) + (\lambda_2\lambda_4+\lambda_3\lambda_5)\cos(\theta-\theta_0)}{\sqrt{\lambda_2^2+\lambda_3^2}}.
\end{align}
\textbf{Геометрия.} Траектория $\vv p(\theta)$ есть эллипс; максимум $p$ может возникать в одной или двух антиподальных точках. Это и порождает одно- и двухимпульсные решения.

\section{Двухимпульсный синтез: \emph{алгебра без скачков}}

\subsection{Уравнение переноса}
Пусть импульсы в $\theta_1,\theta_2$ и их векторы $\Delta\vv v_k = u_k\,\uv s(\theta_k)$ направлены по $\vv p(\theta_k)$:
\(
\uv s(\theta)=\vv p(\theta)/\|\vv p(\theta)\|
\).
Тогда из \eqref{eq:Gmap}
\begin{equation}\label{eq:transfer}
\Delta\vv x = \sqrt{\frac{a_0}{K}}\left[\mathbf{G}(\theta_1)u_1\uv s(\theta_1) + \mathbf{G}(\theta_2)u_2\uv s(\theta_2)\right].
\end{equation}
Необходимость оптимальности даёт равенство максимумов $p(\theta_1)=p(\theta_2)$ и, следовательно, \emph{равные нормы} импульсов: $u_1=u_2=u/2$ (доказательство в Доп.\,B).

\subsection{Параметризация через $(R,C,S)$}
Обозначим $C_k=\cos(\theta_k-\theta_0)$, $S_k=\sin(\theta_k-\theta_0)$ и
\(
C=\cos(\phi),\ S=\sin(\phi)
\)
для среднего/полусуммарного угла (см. рисунок \ref{fig:angles}). Из условия максимумов получаем симметрию:
\(
\theta_2-\theta_0=-(\theta_1-\theta_0)\Rightarrow C_1=C_2=C,\ S_1=-S_2=-S
\).
Подстановка в \eqref{eq:transfer} и явное деление на $u$ приводят к \textbf{открытой} системе (все алгебраические шаги расписаны в Доп.\,C):
\begin{align}
\frac{\Delta (a/a_0)}{u} &= \frac{2C\sqrt{1+R^2}}{D},\label{eq:Da}\\
\frac{\Delta e_y}{u} &= \frac{2(C^2+R^2)}{\sqrt{1+R^2}\,D},\qquad
\frac{\Delta e_x}{u} = \frac{2CS\,(1-2\delta)}{\sqrt{1+R^2}\,D},\label{eq:De}\\
\frac{\Delta i_y}{u} &= \frac{2RS^2}{\sqrt{1+R^2}\,D},\qquad
\frac{\Delta i_x}{u} = \frac{2RCS\,(1-2\delta)}{\sqrt{1+R^2}\,D},\label{eq:Di}
\end{align}
где
\(
D=4R^2+(1-3R^2)C^2
\)
и
\(
\delta=u_2/u\in[0,1]
\)
--- доля второго импульса (её знак ответственен за асимметрию по касательной и нормали).

\subsection{Ликвидация углов и доли импульса}
Введём наблюдаемые углы ориентации
\(
\tan\omega=\Delta e_y/\Delta e_x,\
\tan\Omega=\Delta i_y/\Delta i_x
\),
и их разность $\omega-\Omega$. Делим \eqref{eq:De} на \eqref{eq:Di} (в компонентах ``$x$''), получая связь для $(1-2\delta)$; затем подставляем её в отношение ``$y$''-компонент, получая \textbf{замкнутое уравнение для $R$}. В явном виде (полное выведение в Доп.\,C):
\begin{equation}\label{eq:R_closed}
\boxed{\;
\frac{\Delta a^2/a_0^2+\Delta i^2-\Delta e^2}{\Delta i\,\Delta e \sin(\omega-\Omega)}
= \frac{2R^2+1}{\sqrt{1+R^2}}\,\frac{1}{\sqrt{1-3R^2}}
\;}
\end{equation}
(ветвление по знакам выбрано так, чтобы соблюсти пределы $\Delta i\to 0$ и $\Delta e\to0$).
Решая \eqref{eq:R_closed} относительно $R^2$, получаем \emph{квадратичное} уравнение относительно $z=R^2$, приводящее к компактной формуле
\begin{equation}\label{eq:R_formula}
R^2=1+\frac12\Xi^2-\frac12\sqrt{\Xi^2+\frac14\Xi^4},\qquad
\Xi\equiv \frac{\Delta a^2/a_0^2+\Delta i^2-\Delta e^2}{\Delta i\,\Delta e \sin(\omega-\Omega)}.
\end{equation}
Далее восстанавливаются $C,S$ из \eqref{eq:Da}--\eqref{eq:Di}, а $1-2\delta$ из отношения ``$x$''-компонент (см. Доп.\,C).

\subsection{Суммарный минимум-расход}
Подставляя найденные $R,C,S$ в \eqref{eq:Da}--\eqref{eq:Di} и устраняя их, выводим суммарный оптимальный расход
\begin{equation}\label{eq:u_final}
\boxed{\;
u=\sqrt{\frac{K}{2a_0}}\,
\sqrt{
\Delta i^2+\Delta e_x^2+\Delta e_y^2-\frac{\Delta a^2}{2a_0^2}
+\sqrt{\Bigl(\Delta i^2-\Delta e_x^2-\Delta e_y^2+\frac{\Delta a^2}{a_0^2}\Bigr)^2+4\Delta i^2\,\Delta e_y^2}
}\;.
\end{equation}
}
\noindent
\textbf{Проверка границ.} (i) $\Delta i=0$ даёт две компланарные ветви (пересекающиеся/непересекающиеся); (ii) $\Delta e=0$ даёт $u=\sqrt{K/a_0}\,\Delta i$; (iii) инвариантность при $\Delta e_y\mapsto -\Delta e_y$ реализована лишь через смешанный член под радикалом --- как и предсказывает симметрия модели.

\section{Классификация решений и условия применимости}
\begin{itemize}
\item \textbf{Узловой случай} ($\lambda_1=0$): импульсы через $\pi$, касательная меняет знак, нормаль равна по модулю; область применимости:
$|\Delta a/a_0|\le |\Delta e_x|$ и $\Delta i^2\ge 3\Delta e_y^2$ (вывод из структуры $D$ при $C=0$).
\item \textbf{Невырожденный} ($0<R<1/\sqrt{3}$): общий двухимпульсный случай, параметры $R,C,S,\delta$ восстанавливаются по алгоритму Доп.\,C.
\item \textbf{Сингулярный} ($R=1/\sqrt{3}$): круговой \emph{primer}-локус, отсюда критическое соотношение для $\Delta a$; соответствующие формулы получаем из предела $R\to 1/\sqrt{3}$.
\end{itemize}

\section{Сопряжённые множители как градиент расхода}
Коэффициенты $\lambda_i$ находятся из стационарности \eqref{eq:u_final} по приращениям элементов (детально в Доп.\,D).
Приведём результат в компактной форме, чтобы не загромождать основной поток:
\begin{align}
\lambda_1 &= \frac{K\,\Delta a}{2a_0^2\,u}\left[-\frac12+\frac{\Delta i^2-\Delta e_x^2-\Delta e_y^2+\Delta a^2/a_0^2}{\sqrt{\bigl(\Delta i^2-\Delta e_x^2-\Delta e_y^2+\Delta a^2/a_0^2\bigr)^2+4\Delta i^2\Delta e_y^2}}\right],\\
\lambda_2 &= \frac{K\,\Delta e_y}{2a_0\,u}\left[1-\frac{-\Delta i^2-\Delta e_x^2-\Delta e_y^2+\Delta a^2/a_0^2}{\sqrt{\bigl(\Delta i^2-\Delta e_x^2-\Delta e_y^2+\Delta a^2/a_0^2\bigr)^2+4\Delta i^2\Delta e_y^2}}\right],\\
\lambda_3 &= \frac{K\,\Delta e_x}{2a_0\,u}\left[1-\frac{\Delta i^2-\Delta e_x^2-\Delta e_y^2+\Delta a^2/a_0^2}{\sqrt{\bigl(\Delta i^2-\Delta e_x^2-\Delta e_y^2+\Delta a^2/a_0^2\bigr)^2+4\Delta i^2\Delta e_y^2}}\right],\\
\lambda_4 &= -\frac{K\,\Delta e_x\,\Delta e_y}{a_0\,u}\,\frac{\Delta i}{\sqrt{\bigl(\Delta i^2-\Delta e_x^2-\Delta e_y^2+\Delta a^2/a_0^2\bigr)^2+4\Delta i^2\Delta e_y^2}},\\
\lambda_5 &= \frac{K\,\Delta i}{2a_0\,u}\left[1+\frac{\Delta i^2-\Delta e_x^2+\Delta e_y^2+\Delta a^2/a_0^2}{\sqrt{\bigl(\Delta i^2-\Delta e_x^2-\Delta e_y^2+\Delta a^2/a_0^2\bigr)^2+4\Delta i^2\Delta e_y^2}}\right].
\end{align}

\section{Алгоритм восстановления решения из данных перехода}
\begin{enumerate}[label=\arabic*.]
\item Вычислить $\Delta e=\sqrt{\Delta e_x^2+\Delta e_y^2}$, $\Delta i=\sqrt{\Delta i_x^2+\Delta i_y^2}$, углы $\omega,\Omega$.
\item Найти $R$ по \eqref{eq:R_formula}, затем $C$ из \eqref{eq:Da}, $S$ из \eqref{eq:Di}.
\item Найти $1-2\delta$ из отношения для ``$x$''-компонент \eqref{eq:De}--\eqref{eq:Di}, затем $\delta$.
\item Углы импульсов: $\theta_{1,2}=\theta_0\mp \arctan2(S,C)$, где $\theta_0=\arctan2(\lambda_2,\lambda_3)$, а $(\lambda_2,\lambda_3)$ из градиента $u$.
\item Векторы импульсов: $\Delta\vv v_k=(u/2)\,\vv p(\theta_k)/\|\vv p(\theta_k)\|$ с $\vv p$ из секции primer-вектора.
\end{enumerate}

\section*{Заключение}
Представлен \textbf{полный, линейный, пошаговый} вывод формул с явным устранением промежуточных параметров. Отдельные места статьи, где обычно встречаются опечатки, снабжены проверками пределов и симметрий.

\appendix
\section{Дополнение A: от уравнений Гаусса к \eqref{eq:lin_a}--\eqref{eq:lin_ix}}
Подробная линеаризация: \emph{каждый} шаг с отбрасыванием членов $O(e^2,i^2,ei)$, переход к переменной $u$, а также вывод матрицы $\mathbf{G}(\theta)$ для мгновенного импульса.

\section{Дополнение B: равенство норм импульсов на максимумах $p$}
Доказательство через необходимое условие ПМП: при двух импульсах максимум $p$ должен быть равен в точках, иначе перенос части нормы уменьшает $J$.

\section{Дополнение C: вывод \eqref{eq:R_closed}, формул \eqref{eq:Da}--\eqref{eq:Di} и восстановления $\delta$}
Алгебра без пропусков: деления компонент, исключение $C,S$ и получение уравнения на $R$.

\section{Дополнение D: градиент \eqref{eq:u_final} и формулы для $\lambda_i$}
Аккуратное дифференцирование с цепным правилом и нормировкой $p_{\max}=1$.

\end{document}
